\subsection{Simple non-negative functions}
\lecdate{Lec 20 - Mar 26 (Week 11)}
added end of lusins pf, pf of that other lemma.
wasnt in class for it tho

we start our thy of integration by examining non-negative simple fns.
we then move on to general mble non-neg fns, and finally general mble fns.

\begin{defn}
	let $\phi:E\gto\bR$ be simple w representation:
	\begin{equation*}
		\phi \ = \ \sum_{j=1}^{n}\alpha_{j}\chi_{E_{j}}
	\end{equation*}
	where $0\leq\alpha_{j}<\infty$, $E_{j}\in\cl{E}$ disj.
	we define its \textbf{lebesgue integral} as:
	\begin{equation*}
		\int_{E}\phi \ = \ \sum_{j=1}^{n}\alpha_{j}m(E_{j}) \ \ \in [0,\infty]
	\end{equation*}
	we use the convention that $0\cdot \infty=0$.
\end{defn}
this should be obv, ie, there isnt rly anything else the integral could be.
however, since such reps of simple fns are not necessarily unique,
we should check this is well-defnd.

\begin{lm}
	sps $\phi$ has image $\set{\beta_{1},\dots,\beta_{m}}$, and let $F_{j}=\phi^{-1}(\beta_{j})$.
	then:
	\begin{equation*}
		\int_{E}\phi \ = \ \sum_{j=1}^{m}\beta_{j}m(F_{j})
	\end{equation*}
\end{lm} \

\begin{pf}
	let $\phi=\sum_{j=1}^{n}\alpha_{j}\chi_{E_{j}}$, $E_{j}$ disj. then:
	\begin{align*}
		\int_{E}\phi \ = \ \sum_{j=1}^{n}\alpha_{j}m(E_{j})
		\ = \ \sum_{\substack{j=1 \\ \alpha_{j}\neq0}}^{n}\alpha_{j}m(E_{j})
		& \ = \ \sum_{\substack{k=1 \\ \beta_{k}\neq0}}^{m}\sum_{\substack{j=1 \\ \alpha_{j}=\beta_{k}}}^{n}\alpha_{j}m(E_{j}) \\
		& \ = \ \sum_{\substack{k=1 \\ \beta_{k}\neq0}}^{m}\sum_{\substack{j=1 \\ \alpha_{j}=\beta_{k}}}^{n}\beta_{k}m(E_{j}) \\
		& \ = \ \sum_{\substack{k=1 \\ \beta_{k}\neq0}}^{m}\beta_{k}\sum_{\substack{j=1 \\ \alpha_{j}=\beta_{k}}}^{n}m(E_{j}) \\
		& \ = \ \sum_{\substack{k=1 \\ \beta_{k}\neq0}}^{m}\beta_{k}m(F_{k}) \\
		& \ = \ \sum_{k=1}^{m}\beta_{k}m(F_{k})
	\end{align*}
	note the use of additivity of measure in the 2nd last equality.
\end{pf} \

\begin{prop}
	let $\phi,\psi:E\gto\bR$ be simple non-neg fns, $\lambda\in[0,\infty)$. then:
	\begin{equation*}
		\int_{E}(\phi+\psi) \ = \ \int_{E}\phi+\int_{E}\psi \qquad
		\int_{E}\lambda\phi \ = \ \lambda\int_{E}\phi
	\end{equation*}
	note we implicitly assume that $\phi+\psi$ and $\lambda\phi$ are non-neg simple.
\end{prop} \

\begin{pf}
	scaling is clear.
	let $\phi=\sum_{j=1}^{n}\alpha_{j}\chi_{E_{j}}$ and $\psi=\sum_{k=1}^{m}\beta_{k}\chi_{F_{k}}$,
	where $E_{j}$ disj, $E=\bigcup E_{j}$, similarly for $F_{k}$. then:
	\begin{equation*}
		\phi+\psi \ = \ \sum_{j=1}^{n}\sum_{k=1}^{m}(\alpha_{j}+\beta_{k})\chi_{E_{j}\cap F_{k}}
	\end{equation*} \vspace{-0.4in}
	\begin{align*}
		\ \implies \ \int_{E}\phi+\psi
		& \ = \ \sum_{j=1}^{n}\sum_{k=1}^{m}(\alpha_{j}+\beta_{k})m(E_{j}\cap F_{k}) \\
		& \ = \ \sum_{j=1}^{n}\sum_{k=1}^{m}\alpha_{j}m(E_{j}\cap F_{k})
		+ \sum_{k=1}^{m}\sum_{j=1}^{n}\beta_{k}m(E_{j}\cap F_{k}) \\
		& \ = \ \sum_{j=1}^{n}\alpha_{j}\sum_{k=1}^{m}m(E_{j}\cap F_{k})
		+ \sum_{k=1}^{m}\beta_{k}\sum_{j=1}^{n}m(E_{j}\cap F_{k}) \\
		& \ = \ \sum_{j=1}^{n}\alpha_{j}m(E_{j})+\sum_{k=1}^{m}\beta_{k}m(F_{k}) \\
		& \ = \ \int_{E}\phi+\int_{E}\psi
	\end{align*}
\end{pf} \
we only rly have one last srs pf left, then everythign else will follow nicely.

\begin{crll}
	for any $\phi=\sum_{j=1}^{n}\alpha_{j}\chi_{E_{j}}$ w $\alpha_{j}\in[0,\infty)$,
	$E_{j}\in\cl{E}$ not necessarily disj, we have:
	\begin{equation*}
		\int_{E}\phi \ = \ \sum_{j=1}^{n}\alpha_{j}m(E_{j})
	\end{equation*}
\end{crll} \

\begin{pf}
	by linearity:
	\begin{equation*}
		\int_{E}\phi \ = \ \int_{E}\sum_{j=1}^{n}\alpha_{j}\chi_{E_{j}}
		\ = \ \sum_{j=1}^{n}\int_{E}\alpha_{j}\chi_{E_{j}}
		\ = \ \sum_{j=1}^{n}\alpha_{j}m(E_{j})
	\end{equation*}
\end{pf} \

\begin{crll}
	if $\phi,\psi$ simple, $0\leq\psi\leq\phi$, then:
	\begin{equation*}
		\int_{E}\psi \ \leq \ \int_{E}\phi
	\end{equation*}
\end{crll} \

\begin{pf}
	note $\phi=(\phi-\psi)+\psi$, and $\phi-\psi$ is simple non-neg by assumption.
	by linearity:
	\begin{equation*}
		\int_{E}\phi \ = \ \int_{E}\phi-\psi + \int_{E}\psi
		\ \geq \ \int_{E}\psi
	\end{equation*}
\end{pf}
